\section{Background and Motivation}
\label{sec:background}

\subsection{Python Native Extensions}

CPython exposes a C interface for defining modules, types, and functions and
for manipulating interpreter objects. A native function commonly parses
Python arguments, creates or accesses \code{PyObject} values, calls other
Python/C APIs, and returns either a Python object or an error indicator. The
interface uses conventions rather than a C type system to encode many safety
properties. A \emph{new reference} transfers ownership to the caller, while a
\emph{borrowed reference} remains owned elsewhere. Some APIs steal ownership
of an argument. Most object-producing APIs return \code{NULL} on failure and
set the thread-local exception indicator. Other APIs use integer sentinels,
and several resource APIs require a matching release operation
\cite{pythoncapi}.

The conventions interact. If \code{PyTuple_New} returns \code{NULL}, native
code must avoid dereferencing the result, preserve the pending exception, and
release objects acquired earlier on that path. If \code{PyTuple_SetItem}
succeeds, it steals an item reference, so a later decrement can access a freed
object. Correctness therefore depends on the API contract, the program object
to which the contract applies, and the path between acquisition and exit.

\subsection{Python/C API Safety Requirements}

We call a requirement \emph{API-induced} when a Python/C API call or a
Python-runtime callback creates the relevant duty. Reference ownership is the
most visible example, but the same structure appears in other parts of the
interface. A successful \code{PyObject\_GetBuffer} call creates a duty to call
\code{PyBuffer\_Release}; a fallible conversion creates a duty to check its
sentinel before use; releasing the global interpreter lock creates a duty to
reacquire it before another Python/C API call; and a garbage-collected type
creates a multi-part lifecycle duty for allocation, tracking, traversal,
clearing, and deallocation.

Three properties make these requirements unsuitable for simple call-pair
matching. A requirement can admit several correct outcomes, such as release,
transfer, or return. The appropriate outcome can differ across paths. In
addition, a syntactically present action can be semantically wrong, as when a
cleanup block decrements a borrowed reference or frees memory with an
allocator from another domain. These properties motivate an explicit
representation of both the duty and its admissible discharge conditions.

\subsection{Motivating Example}

Figure~\ref{fig:motivation} presents a defect found in the SciPy
\code{fitpack\_sphere} wrapper. The function receives \code{tt} and \code{tp}
from \code{PyArg\_ParseTuple}. When \code{iopt == -1}, it assigns these
borrowed references to \code{ap\_tt} and \code{ap\_tp}. A later validation
failure reaches the shared cleanup block, which decrements both variables.
The cleanup is correct when \code{iopt == 0} because that branch creates new
arrays, but it is unsafe for the borrowed-reference branch. The decrement can
prematurely release an object still owned by the caller.

\begin{figure}[t]
\begin{lstlisting}[language=C,numbers=left,firstnumber=1]
if (iopt == -1) {
    ap_tt = tt;                 // borrowed
    ap_tp = tp;                 // borrowed
} else if (iopt == 0) {
    ap_tt = PyArray_SimpleNew(...); // new
    if (ap_tt == NULL) goto fail;
    ap_tp = PyArray_SimpleNew(...); // new
    if (ap_tp == NULL) goto fail;
}
if (invalid_dimensions) goto fail;
...
fail:
Py_XDECREF(ap_tt);
Py_XDECREF(ap_tp);
return NULL;
\end{lstlisting}
\caption{Path-dependent ownership in \code{fitpack\_sphere}. The cleanup
releases owned arrays created by the second branch but also releases borrowed
arrays assigned by the first branch.}
\label{fig:motivation}
\end{figure}

A source-level pattern such as ``assignment followed by decrement'' cannot
distinguish the branches. Tracking only the final reference-count delta is
also insufficient unless the analysis recovers the ownership semantics of
both assignments and the feasibility of each cleanup path. Under our
formulation, the two assignments bind a reference-ownership obligation. The
new-reference branch is discharged by cleanup on failure and ownership
transfer on success. The borrowed-reference branch has no owned reference to
release, and the decrement violates the post-acquisition condition. The
example illustrates why the checker must combine API semantics with
path-specific discharge reasoning.
